\documentclass[12pt]{article}
\usepackage{fullpage}
\usepackage{amsmath,amssymb,amsthm}

\newtheorem{lemma}{Lemma}
\newtheorem{proposition}{Proposition}
\newtheorem{theorem}{Theorem}

\DeclareMathOperator{\Lk}{Lk}
\newcommand{\bbarotimes}{\mathbin{\overline{\otimes}}}
\newcommand{\cA}{\mathcal A}
\newcommand{\cB}{\mathcal B}
\newcommand{\cH}{\mathcal H}
\newcommand{\K}{\mathbb K}
\newcommand{\C}{\mathbb C}

\begin{document}

\title{Proper proximality of irreducible graph-product von Neumann algebras}
\author{}
\date{}
\maketitle

\section*{Problem statement and interpretation}

I use the standard graph-product meaning of irreducible: the finite graph
$\Gamma=(V,E)$ is not a non-trivial join, equivalently the complement graph
$\Gamma^c$ is connected.  Traces are faithful normal tracial states.  For
$U\subset V$, $M_U$ denotes the graph product over the induced subgraph and
$E_U:M\to M_U$ the canonical trace-preserving expectation.

Let \(A_N=\mathbb B(L^2N)\) and
\[
        \cA_N=(A_N^{\sharp_J})^*
        =p_{\rm nor}A_N^{**}p_{\rm nor}
\]
be the normal bidual of Ding--Kunnawalkam Elayavalli--Peterson
\cite{DKEP}.  If \(X\) is an \(N\)-boundary piece, \(K_X(N)\) is DKEP's
hereditary algebra associated with \(X\), and \(q_X\in\cA_N\) is the support of
\((K_X(N)^{\sharp_J})^*\).  Thus
\[
 \widetilde S_X(N)=
 \{T\in\cA_N:[T,a]\in q_X\cA_Nq_X\hbox{ for all }a\in JNJ\}.
\]
For \(X=\mathbb K(L^2N)\) write \(q_0\) and \(\widetilde S(N)\).  DKEP
Lemma 8.5 says that, for a separable finite factor, proper proximality relative
to \(X\) is equivalent to the non-existence of an \(N\)-central state on
\(\widetilde S_X(N)\) restricting to the trace.

\section*{Boundary support calculus}

We record the two pieces of support calculus used below.  First, if
\(B\subset A_N\) is a \(C^*\)-algebra which is an \(N\)- and \(JNJ\)-bimodule
with faithful multiplier actions, and if \(s_B\in A_N^{**}\) is its ordinary
open support, DKEP Proposition 2.10 gives
\[
        (B^{\sharp_J})^*=p_{\rm nor}s_BA_N^{**}s_Bp_{\rm nor},       \tag{1}
\]
and \(s_B\) commutes with \(p_{\rm nor}\).  In particular \(q_X=p_{\rm nor}s\),
where \(s\) is the ordinary support of \(K_X(N)\).  DKEP Proposition 3.6
identifies \(K_X(N)^{\infty,1}\) with the \(\|\cdot\|_{\infty,1}\)-closure of
the original boundary piece \(X\), so this is precisely the normal-bidual
support used in Lemma 8.5.

Second, let \(B,C\subset A_N\) have ordinary supports \(r,s\).  If \(e\) is a
projection in \(A_N^{**}\) and
\[
        BC\subset eA_N^{**}s,\qquad CB\subset sA_N^{**}e ,
\]
then
\[
        r(sA_N^{**}s)r\subset eA_N^{**}e .                         \tag{2}
\]
Indeed, take approximate units \(b_i\to r\), \(c_j\to s\) weak-*, and pass to
successive weak-* limits in \(b_ic_jxc_{j'}b_{i'}\in eA_N^{**}e\).

We also use the following elementary rectangular version.  If \(D\subset A_N\)
has support \(d\) and \(T:L^2N\to H\) is bounded with \(T^*T\in D\), then
\(T=T d\) in the bidual sense: for an approximate unit \(u_i\) of \(D\),
\(Tu_i\to T\) strongly because
\(\|T(1-u_i)\xi\|^2=\langle (1-u_i)T^*T(1-u_i)\xi,\xi\rangle\to0\).
Consequently, if also \(S^*S\in C\) with support \(s\), then
\[
        T^*S\in dA_N^{**}s .                                      \tag{3}
\]

For \(U\subset V\), put \(\cH_U=L^2M\otimes_{M_U}L^2M\).  The vector
\(\hat1\otimes\hat1\) has coefficient projection \(P_U=e_{M_U}\), so DKEP
Example 3.4 realizes \(X_{M_U}\) as the hereditary algebra generated by the
\(M,JMJ\)-translates of this coefficient.  From the definitions
\(K_X^L=\overline{A_MX}^{\,\|\cdot\|_{\infty,2}}\) and
\(K_X=(K_X^L)^*K_X^L\), and from
\(T_{\bar\xi}=T_\xi^*\), \(T_{\xi\otimes\eta}=T_\eta T_\xi\), it follows that
\(K_{X_{M_U}}(M)\) is the closed span of coefficients obtained from non-empty
finite fusions of \(\cH_U\) and \(\overline{\cH_U}\).  CDHJKN's fusion rule then
rewrites these fusions as direct sums of \(\cH_Z\)'s with \(Z\subset U\).
Accordingly let \(\cB_U\) be the norm-closed span of all coefficients
\(T_\xi^*T_\eta\), where \(\xi,\eta\) are bounded vectors in a non-empty finite
fusion of correspondences \(\cH_Z\) and \(\overline{\cH_Z}\) with \(Z\subset U\);
then \(\cB_U=K_{X_{M_U}}(M)\).  It is a \(C^*\)-algebra since products and
adjoints are again coefficients:
\[
 (T_\xi^*T_\eta)(T_\alpha^*T_\beta)
   =T_{\xi\otimes\bar\eta\otimes\alpha}^{\,*}T_\beta,\qquad
 (T_\xi^*T_\eta)^*=T_\eta^*T_\xi .
\]
Let \(s_U\) be the ordinary support of \(\cB_U\).  By (1),
\[
        q_U:=q_{X_{M_U}}=p_{\rm nor}s_U .                          \tag{4}
\]

\begin{lemma}[intersection of graph-product boundary corners]\label{lem:intersection}
For all \(U,W\subset V\),
\[
        q_W(q_U\cA_Mq_U)q_W\subset q_{U\cap W}\cA_Mq_{U\cap W}.     \tag{5}
\]
Consequently, if \(U_1,\ldots,U_m\subset V\) and
\[
R_m=q_{U_m}q_{U_{m-1}}\cdots q_{U_1}q_{U_2}\cdots q_{U_{m-1}}q_{U_m},
\]
then \(R_m\in q_I\cA_Mq_I\), where \(I=\cap_iU_i\).
\end{lemma}

\begin{proof}
Put \(I=U\cap W\).  We first prove, in \(A_M^{**}\),
\[
        \cB_W\cB_U\subset s_IA_M^{**}s_U,\qquad
        \cB_U\cB_W\subset s_UA_M^{**}s_I .                         \tag{6}
\]
It suffices by norm closure to multiply two coefficient operators.  If the
first coefficient comes from a finite fusion using subsets of \(W\) and the
second from one using subsets of \(U\), then
\[
 (T_{\xi_1}^*T_{\xi_2})(T_{\eta_1}^*T_{\eta_2})
 =
 T_{\xi_1\otimes\overline{\xi_2}\otimes\eta_1}^{\,*}T_{\eta_2}.    \tag{7}
\]
The CDHJKN bimodule decomposition and fusion rule \cite[Theorem 5.4 and
Proposition 5.5]{CDHJKN} say that every such finite fusion over subsets of
\(W\) decomposes as a direct sum of copies of \(\cH_L\) with \(L\subset W\),
and similarly for \(U\); after fusing the two \(W\)-parts with the \(U\)-part,
all summands containing
\(\xi_1\otimes\overline{\xi_2}\otimes\eta_1\) are \(\cH_L\)'s with
\(L\subset I\).  Finite-support vectors in these direct sums are dense, and the
coefficient maps are norm continuous, so the assertion holds for arbitrary
bounded vectors.  Thus
\(T_{\alpha}^*T_{\alpha}\in\cB_I\) for
\(\alpha=\xi_1\otimes\overline{\xi_2}\otimes\eta_1\), while
\(T_{\eta_2}^*T_{\eta_2}\in\cB_U\).  The rectangular support lemma (3) gives
\(T_\alpha^*T_{\eta_2}\in s_IA_M^{**}s_U\).  This proves the first inclusion
in (6), and the second follows by adjoints.

Apply (2) with \(B=\cB_W\), \(C=\cB_U\), \(e=s_I\).  We obtain
\[
        s_W(s_UA_M^{**}s_U)s_W\subset s_IA_M^{**}s_I .             \tag{8}
\]
Since the \(s_\bullet\)'s commute with \(p_{\rm nor}\), cutting (8) by
\(p_{\rm nor}\) and using (4) gives (5).  The palindromic statement follows by
induction from
\(R_{k+1}=q_{U_{k+1}}R_kq_{U_{k+1}}\) and (5).
\end{proof}

\begin{lemma}[finite palindromic criterion]\label{lem:finite}
Let \(N\) be a separable non-amenable finite factor.  Let \(X_1,\ldots,X_m\)
be boundary pieces such that, with \(q_i=q_{X_i}\),
\[
        R=q_mq_{m-1}\cdots q_1q_2\cdots q_{m-1}q_m\in q_0\cA_Nq_0. \tag{9}
\]
If \(N\) is properly proximal relative to each \(X_i\), then \(N\) is properly
proximal.
\end{lemma}

\begin{proof}
Assume \(N\) is not properly proximal.  DKEP Lemma 8.5 gives an \(N\)-central
state \(\omega\) on \(\widetilde S(N)\) with \(\omega|_N=\tau\).  Fix \(i\) and
write \(q=q_i\).  Since \(q\) commutes with \(N\) and \(JNJ\), for
\(T\in\widetilde S_{X_i}(N)\) one has
\(q^\perp Tq^\perp\in\widetilde S(N)\).  If \(\omega(q^\perp)>0\), the
compression \(T\mapsto \omega(q^\perp Tq^\perp)/\omega(q^\perp)\) is an
\(N\)-central state on \(\widetilde S_{X_i}(N)\).  Its restriction to \(N\) is
normal and central; because \(N\) is a factor it is \(\tau\), contradicting
relative proper proximality.  Hence \(\omega(q_i^\perp)=0\) for every \(i\).

Extend \(\omega\) to a state \(\bar\omega\) on \(\cA_N\).  In the GNS space of
\(\bar\omega\), each \(q_i\) fixes the cyclic vector, and therefore
\(\bar\omega(R)=1\).  Since \(0\le R\le q_0\), \(\bar\omega(q_0)=1\).  The
natural DKEP normal-bidual map
\[
        \Theta:\mathbb B(L^2N)\longrightarrow q_0\cA_Nq_0
\]
is \(N\)-bimodular, u.c.p., and agrees with the identity on \(N\) after the
canonical embedding.  Thus \(\bar\omega\circ\Theta\) is an \(N\)-central state
on \(\mathbb B(L^2N)\) extending \(\tau\), i.e. a Connes hypertrace.  This
contradicts non-amenability of the finite factor \(N\).
\end{proof}

\section*{A relative Bass--Serre paradox}

\begin{proposition}\label{prop:BS}
Let \(N=P*_{B}Q\) with trace-preserving expectations, and suppose
\(Q=B\bbarotimes A\).  Assume that \(A\) contains a trace-zero unitary \(a\) and
that \(P\) contains unitaries \(w_1,w_2\) satisfying
\[
        E_B(w_i)=0,\qquad E_B(w_1^*w_2)=0 .
\]
Then \(N\) is properly proximal relative to \(X_B\).
\end{proposition}

\begin{proof}
Write \(H_P=L^2P\ominus L^2B\), \(H_Q=L^2Q\ominus L^2B\), and use the usual
amalgamated-free-product decomposition of \(L^2N\).  Let \(p\) and \(q\) be
the projections onto the sums of reduced tensors whose first letter is in
\(H_P\) and \(H_Q\), respectively; then \(1=e_B+p+q\).

Let \(R_x=Jx^*J\).  On reduced tensors, if \(x\in P\ominus B\), right
multiplication by \(x\) preserves the first letter except on \(L^2B\) and on
one-letter \(P\)-words.  Hence
\[
        [p,R_x]=R_xe_B-e_BR_x,\qquad [q,R_x]=0 .
\]
For \(x\in Q\ominus B\) the analogous identities hold with \(p\) and \(q\)
interchanged, and elements of \(B\) commute with both projections.  The two
commutators displayed lie in \(X_B\), since \(X_B\) contains the left and right
\(N\)-translates of \(e_B\).  DKEP Lemma 6.1 therefore gives
\(p,q\in S_{X_B}(N)\).

Let \(\varphi\) be an \(N\)-central state on \(S_{X_B}(N)\) normal on \(N\).
Left multiplication by \(w_i\) sends \(qL^2N\) into \(pL^2N\).  The two images
are orthogonal: for reduced \(Q\)-initial vectors \(\xi,\eta\), the
\(B\)-valued inner product of \(w_1\xi\) and \(w_2\eta\) contains the centered
\(P\)-letter \(w_1^*w_2-E_B(w_1^*w_2)=w_1^*w_2\).  Thus
\(w_1qw_1^*+w_2qw_2^*\le p\).  Similarly, since \(a\in Q\ominus B\),
\(apa^*\le q\).  Centrality gives
\(2\varphi(q)\le\varphi(p)\le\varphi(q)\), so \(\varphi(p)=\varphi(q)=0\).

Finally \(h=w_1a\) is a \(B\)-Haar unitary.  For \(n>0\), \(h^n\) is the
reduced alternating word \(w_1aw_1a\cdots w_1a\); for \(n<0\) it is the
reduced alternating word \(a^*w_1^*\cdots a^*w_1^*\).  Hence \(E_B(h^n)=0\) for
\(n\ne0\).  Therefore the projections \(h^ne_Bh^{-n}\) are pairwise
orthogonal, because \(e_Bh^{k}e_B=E_B(h^k)e_B\).  They all have the same
\(\varphi\)-value, so \(\varphi(e_B)=0\).  This contradicts
\(\varphi(1)=\varphi(e_B+p+q)=1\).  By DKEP Theorem 6.2, \(N\) is properly
proximal relative to \(X_B\).
\end{proof}

\section*{Application to graph products}

Fix \(v\in V\) and put \(L_v=\Lk_\Gamma(v)\).  The graph product splits as
\[
        M\cong M_{V\setminus\{v\}}*_{M_{L_v}}
        (M_{L_v}\bbarotimes M_v).                                \tag{10}
\]
Let \(C(v)=\{r\ne v:r\not\sim_\Gamma v\}\).  Since \(\Gamma^c\) is connected,
\(C(v)\ne\emptyset\).  If \(C(v)\) contains distinct \(r,t\), choose trace-zero
unitaries \(u_r\in M_r\), \(u_t\in M_t\) and put \(w_1=u_r,w_2=u_t\).  If
\(C(v)=\{r\}\), connectedness of \(\Gamma^c\) and \(|V|\ge3\) give
\(s\in L_v\) with \(s\not\sim_\Gamma r\); put \(w_1=u_r,w_2=u_su_r\).  In both
cases reduced-word freeness gives
\[
        E_{M_{L_v}}(w_i)=0,\qquad
        E_{M_{L_v}}(w_1^*w_2)=0 .
\]
The trace-zero unitary in \(M_v\) supplies the unitary \(a\) in the tensor
factor \(M_v\) of (10).  Proposition \ref{prop:BS} shows that \(M\) is properly
proximal relative to \(X_{M_{L_v}}\) for every \(v\in V\).

Enumerate \(V=\{v_1,\ldots,v_m\}\).  Since no vertex belongs to its own link,
\(\cap_i L_{v_i}=\emptyset\).  Lemma \ref{lem:intersection} gives
\[
 q_{L_{v_m}}\cdots q_{L_{v_1}}q_{L_{v_2}}\cdots q_{L_{v_m}}
 \in q_{\emptyset}\cA_Mq_{\emptyset}.
\]
Here \(M_\emptyset=\C\), and \(X_{M_\emptyset}\) is the compact boundary piece,
because \(e_{M_\emptyset}\) is rank one and its left and right \(M\)-translates
generate the finite-rank operators.

It remains only to check the factor and non-amenability hypotheses of Lemma
\ref{lem:finite}.  The graph product is separable.  By CDHJKN Main Theorem
0.5, factoriality can fail only through a universal vertex, or through a
non-edge pair adjacent to all other vertices with both vertex algebras
two-dimensional.  These are respectively an isolated vertex or an isolated edge
in \(\Gamma^c\), impossible because \(\Gamma^c\) is connected and has at least
three vertices.  Thus \(M\) is a factor.  By CDHJKN Proposition 6.3, amenability
would force every non-edge pair to be such an isolated two-dimensional pair,
again impossible in a connected complement graph with at least three vertices.
Hence \(M\) is non-amenable.  Lemma \ref{lem:finite} applies, and \(M\) is
properly proximal.

\begin{thebibliography}{9}

\bibitem{DKEP}
C.~Ding, S.~Kunnawalkam Elayavalli, and J.~Peterson,
\emph{Properly proximal von Neumann algebras},
Duke Math. J. \textbf{172} (2023), 2821--2894.

\bibitem{CDHJKN}
I.~Charlesworth, R.~de Santiago, B.~Hayes, D.~Jekel,
S.~Kunnawalkam Elayavalli, and B.~Nelson,
\emph{On the structure of graph product von Neumann algebras},
Publ. Res. Inst. Math. Sci. \textbf{61} (2025), 713--762.

\end{thebibliography}

\end{document}
